%---------------------------------------------------------------------
% UNIT 3 --- EQUATIONS, INEQUALITIES, ABSOLUTE VALUES, PLANE GEOMETRY
% Source: Section 3.9 Self Assessment Questions, pp. 64-65.
%---------------------------------------------------------------------
\chapter{Equations and Inequalities}

\srcnote{Source: \emph{Business Mathematics} 5405/1429, \S3.9 \saq{Self Assessment Questions}, pp.~64--65.}

\begin{keyidea}[Three tools for the whole unit]
\begin{align*}
\text{Quadratic formula}&: && x=\frac{-b\pm\sqrt{b^{2}-4ac}}{2a}
   \quad\text{for } ax^{2}+bx+c=0\\[2pt]
\text{Absolute value}&: && |u|=|v| \iff u=v \ \text{ or } \ u=-v\\[2pt]
\text{Distance / midpoint}&: &&
   d=\sqrt{(x_2-x_1)^{2}+(y_2-y_1)^{2}},\quad
   M=\left(\tfrac{x_1+x_2}{2},\tfrac{y_1+y_2}{2}\right)
\end{align*}
And one rule that costs more marks than any formula: \textbf{multiplying or
dividing an inequality by a negative number reverses the inequality sign.}
\end{keyidea}

%---------------------------------------------------------------------
\section{First-degree equations}

\begin{question}{Q.1}
Solve the following first-degree equations.
\textbf{(a)} $8x-6=5x+3$ \quad
\textbf{(b)} $-15+35x=8x-9$ \quad
\textbf{(c)} $(x+9)-(-6+4x)+4=0$ \quad
\textbf{(d)} $\dfrac{y}{8}-10=\dfrac{y}{4}-9$
\end{question}

\begin{solution}
\textbf{(a)} Collect $x$ on the left, constants on the right:
\[
8x-5x=3+6 \;\Longrightarrow\; 3x=9 \;\Longrightarrow\; x=3 .
\]
\vcheck{$8(3)-6=18$ and $5(3)+3=18$.}

\medskip
\textbf{(b)}
\[
35x-8x=-9+15 \;\Longrightarrow\; 27x=6 \;\Longrightarrow\;
x=\frac{6}{27}=\frac{2}{9}.
\]
\vcheck{LHS $=-15+35(\tfrac29)=-15+\tfrac{70}{9}=-\tfrac{65}{9}$;
RHS $=8(\tfrac29)-9=\tfrac{16}{9}-9=-\tfrac{65}{9}$.}

\medskip
\textbf{(c)} Clear the brackets, watching the sign on the second one:
\[
(x+9)-(-6+4x)+4 = x+9+6-4x+4 = -3x+19 .
\]
\[
-3x+19=0 \;\Longrightarrow\; x=\frac{19}{3}.
\]

\medskip
\textbf{(d)} Multiply through by the LCD, 8, before doing anything else:
\[
y-80=2y-72 \;\Longrightarrow\; -80+72=2y-y \;\Longrightarrow\; y=-8 .
\]
\vcheck{LHS $=\tfrac{-8}{8}-10=-11$; RHS $=\tfrac{-8}{4}-9=-11$.}

\ans{(a) $x=3$ \quad (b) $x=\tfrac{2}{9}$ \quad (c) $x=\tfrac{19}{3}$ \quad
(d) $y=-8$}
\end{solution}

\begin{pitfall}
Part (c): $-(-6+4x)$ is $+6-4x$, not $-6-4x$. The minus sign distributes over
\emph{both} terms in the bracket.
\end{pitfall}

%---------------------------------------------------------------------
\section{Quadratics by factorisation}

\begin{question}{Q.2}
Solve the following quadratic equations using the factorisation method.
\textbf{(a)} $t^{2}+4t=21$ \quad \textbf{(b)} $4z^{2}+18z-10=0$
\end{question}

\begin{solution}
\textbf{(a)} Move everything to one side first --- factorisation only works
against zero:
\[
t^{2}+4t-21=0 .
\]
Two numbers multiplying to $-21$ and adding to $+4$: these are $+7$ and $-3$.
\[
(t+7)(t-3)=0 \;\Longrightarrow\; t=-7 \ \text{ or } \ t=3 .
\]
\vcheck{$(-7)^{2}+4(-7)=49-28=21$; \ $3^{2}+4(3)=9+12=21$.}

\medskip
\textbf{(b)} Take out the common factor 2 to keep the numbers small:
\[
4z^{2}+18z-10=0 \;\Longrightarrow\; 2\left(2z^{2}+9z-5\right)=0
\;\Longrightarrow\; 2z^{2}+9z-5=0 .
\]
Split the middle term: $2\times(-5)=-10$, and $+10$ and $-1$ give $+9$:
\[
2z^{2}+10z-z-5=0 \;\Longrightarrow\; 2z(z+5)-1(z+5)=0
\;\Longrightarrow\; (2z-1)(z+5)=0 .
\]
\[
z=\tfrac{1}{2} \ \text{ or } \ z=-5 .
\]
\vcheck{$4(\tfrac14)+18(\tfrac12)-10=1+9-10=0$; \
$4(25)+18(-5)-10=100-90-10=0$.}

\ans{(a) $t=-7,\;3$ \quad (b) $z=\tfrac12,\;-5$}
\end{solution}

%---------------------------------------------------------------------
\section{Quadratics by formula}

\begin{question}{Q.3}
Solve using the quadratic formula.
\textbf{(a)} $4x^{2}+3x-1=0$ \quad \textbf{(b)} $4t^{2}-64=0$
\end{question}

\begin{solution}
\textbf{(a)} $a=4$, $b=3$, $c=-1$. Compute the discriminant first --- it tells
you what kind of answer to expect:
\[
b^{2}-4ac=3^{2}-4(4)(-1)=9+16=25>0 \quad\text{(two distinct real roots)}.
\]
\[
x=\frac{-3\pm\sqrt{25}}{2(4)}=\frac{-3\pm 5}{8}
\;\Longrightarrow\;
x=\frac{2}{8}=\frac14 \quad\text{or}\quad x=\frac{-8}{8}=-1 .
\]
\vcheck{$4(\tfrac{1}{16})+3(\tfrac14)-1=\tfrac14+\tfrac34-1=0$; \
$4(1)+3(-1)-1=0$.}

\medskip
\textbf{(b)} $a=4$, $b=0$, $c=-64$:
\[
t=\frac{0\pm\sqrt{0-4(4)(-64)}}{8}=\frac{\pm\sqrt{1024}}{8}=\frac{\pm 32}{8}
=\pm 4 .
\]
Faster without the formula: $4t^{2}=64 \Rightarrow t^{2}=16 \Rightarrow t=\pm4$.

\ans{(a) $x=\tfrac14,\;-1$ \quad (b) $t=\pm 4$}
\end{solution}

\begin{pitfall}
In (b) the answer is $t=\pm 4$, not $t=4$. Taking a square root of both sides
of $t^{2}=16$ produces two roots; dropping the negative one loses half the
marks.
\end{pitfall}

%---------------------------------------------------------------------
\section{Completing the square}

\begin{question}{Q.4}
Solve using the completing-the-square method.
\textbf{(a)} $2y^{2}+5y-3=0$ \quad \textbf{(b)} $s^{2}+s-3=0$
\end{question}

\begin{solution}
\textbf{(a)} Divide through by the leading coefficient so that $y^{2}$ stands
alone, then move the constant across:
\[
y^{2}+\frac{5}{2}y-\frac{3}{2}=0
\;\Longrightarrow\;
y^{2}+\frac{5}{2}y=\frac{3}{2}.
\]
Half the coefficient of $y$ is $\tfrac54$; add its square $\tfrac{25}{16}$ to
both sides:
\[
y^{2}+\frac{5}{2}y+\frac{25}{16}=\frac{3}{2}+\frac{25}{16}
=\frac{24}{16}+\frac{25}{16}=\frac{49}{16},
\]
\[
\left(y+\frac54\right)^{2}=\frac{49}{16}
\;\Longrightarrow\;
y+\frac54=\pm\frac74
\;\Longrightarrow\;
y=\frac{-5\pm 7}{4}.
\]
\[
y=\frac{2}{4}=\frac12 \quad\text{or}\quad y=\frac{-12}{4}=-3 .
\]
\vcheck{$2(\tfrac14)+5(\tfrac12)-3=\tfrac12+\tfrac52-3=0$; \
$2(9)+5(-3)-3=18-15-3=0$.}

\medskip
\textbf{(b)} The leading coefficient is already 1:
\[
s^{2}+s=3 .
\]
Half of 1 is $\tfrac12$; add $\tfrac14$:
\[
s^{2}+s+\frac14=3+\frac14=\frac{13}{4}
\;\Longrightarrow\;
\left(s+\frac12\right)^{2}=\frac{13}{4}
\;\Longrightarrow\;
s+\frac12=\pm\frac{\sqrt{13}}{2},
\]
\[
s=\frac{-1\pm\sqrt{13}}{2}\approx 1.303 \ \text{ or } \ -2.303 .
\]

\ans{(a) $y=\tfrac12,\;-3$ \quad (b) $s=\dfrac{-1\pm\sqrt{13}}{2}$}
\end{solution}

%---------------------------------------------------------------------
\section{Linear inequalities}

\begin{question}{Q.5}
Solve the following inequalities.
\textbf{(a)} $z+6\le 10-z$ \quad \textbf{(b)} $3y+6\ge 3y-5$
\end{question}

\begin{solution}
\textbf{(a)}
\[
z+z\le 10-6 \;\Longrightarrow\; 2z\le 4 \;\Longrightarrow\; z\le 2 .
\]
Dividing by $+2$, so the sign does not turn. Solution set
$(-\infty,\,2]$.

\medskip
\textbf{(b)} Subtract $3y$ from both sides and the variable disappears:
\[
3y+6\ge 3y-5 \;\Longrightarrow\; 6\ge -5 .
\]
This is a true statement containing no $y$ at all. It is true whatever $y$ is,
so \emph{every} real number satisfies the original inequality.
\[
\text{Solution set }=\mathbb{R}=(-\infty,\infty).
\]

\ans{(a) $z\le 2$, i.e.\ $(-\infty,2]$ \quad
(b) all real $y$; the inequality reduces to $6\ge -5$, which is always true.}
\end{solution}

\begin{pitfall}
Part (b) is a deliberate trap. ``No $y$ left'' does not mean ``no solution''.
If what remains is \emph{true} ($6\ge-5$) every $y$ works; if what remains is
\emph{false} (say $6\le-5$) no $y$ works. Read the surviving statement.
\end{pitfall}

%---------------------------------------------------------------------
\section{Quadratic inequalities}

\begin{question}{Q.6}
Solve the following second-degree inequalities.
\textbf{(a)} $6t^{2}+t-12>0$ \quad \textbf{(b)} $2s^{2}-3s-2<0$
\end{question}

\begin{solution}
Method for both: find the roots, then decide which side of them the parabola
lies on the required side of the axis.

\medskip
\textbf{(a)} Roots of $6t^{2}+t-12=0$: discriminant $=1+288=289$, $\sqrt{289}=17$,
\[
t=\frac{-1\pm 17}{12}
\;\Longrightarrow\;
t=\frac{16}{12}=\frac43 \quad\text{or}\quad t=\frac{-18}{12}=-\frac32 .
\]
The coefficient of $t^{2}$ is $6>0$, so the parabola opens \emph{upwards}: it is
below the axis between the roots and above the axis outside them. We want
$>0$, i.e.\ above:
\[
t<-\frac32 \quad\text{or}\quad t>\frac43 .
\]
\vcheck{Test $t=0$ (between the roots): $6(0)+0-12=-12$, which is not $>0$ ---
correctly excluded. Test $t=2$: $24+2-12=14>0$ --- correctly included.}

\medskip
\textbf{(b)} Roots of $2s^{2}-3s-2=0$: discriminant $=9+16=25$,
\[
s=\frac{3\pm 5}{4} \;\Longrightarrow\; s=2 \quad\text{or}\quad s=-\frac12 .
\]
Coefficient of $s^{2}$ is $2>0$, so the parabola is below the axis
\emph{between} the roots. We want $<0$:
\[
-\frac12<s<2 .
\]
\vcheck{Test $s=0$: $-2<0$ --- included. Test $s=3$: $18-9-2=7$, not $<0$ ---
excluded.}

\ans{(a) $t<-\tfrac32$ or $t>\tfrac43$ \quad
(b) $-\tfrac12<s<2$}
\end{solution}

%---------------------------------------------------------------------
\section{Absolute value equations and inequalities}

\begin{question}{Q.7}
Solve the following.
\textbf{(a)} $|2y+5|=|y-4|$ \quad
\textbf{(b)} $|y|=|-y+7|$ \quad
\textbf{(c)} $|z^{2}-8|\le 8$ \quad
\textbf{(d)} $|z^{2}-2|\ge 2$
\end{question}

\begin{solution}
\textbf{(a)} $|u|=|v|$ splits into two cases, $u=v$ and $u=-v$:

\emph{Case 1:} $2y+5=y-4 \Rightarrow y=-9$.

\emph{Case 2:} $2y+5=-(y-4)=-y+4 \Rightarrow 3y=-1 \Rightarrow y=-\tfrac13$.

\vcheck{$y=-9$: $|-18+5|=13$ and $|-9-4|=13$. \
$y=-\tfrac13$: $|-\tfrac23+5|=\tfrac{13}{3}$ and
$|-\tfrac13-4|=\tfrac{13}{3}$.}

\medskip
\textbf{(b)} \emph{Case 1:} $y=-y+7 \Rightarrow 2y=7 \Rightarrow y=\tfrac72$.

\emph{Case 2:} $y=-(-y+7)=y-7 \Rightarrow 0=-7$, a contradiction, so this case
yields nothing.

Only $y=\tfrac72$.

\medskip
\textbf{(c)} $|u|\le a$ (with $a\ge0$) is the double inequality $-a\le u\le a$:
\[
-8\le z^{2}-8\le 8
\;\Longrightarrow\;
0\le z^{2}\le 16 .
\]
The left half, $z^{2}\ge 0$, is automatically true for every real $z$ and
imposes nothing. The right half gives
\[
z^{2}\le 16 \;\Longrightarrow\; -4\le z\le 4 .
\]

\medskip
\textbf{(d)} $|u|\ge a$ splits \emph{outwards}: $u\ge a$ or $u\le -a$:
\[
z^{2}-2\ge 2 \quad\text{or}\quad z^{2}-2\le -2,
\]
\[
z^{2}\ge 4 \quad\text{or}\quad z^{2}\le 0 .
\]
The first gives $z\le -2$ or $z\ge 2$. The second, $z^{2}\le 0$, holds only at
the single point $z=0$ --- easy to discard by reflex, but it is a genuine part
of the solution set.
\[
\text{Solution: } z\le -2,\quad z=0,\quad\text{or}\quad z\ge 2 .
\]
\vcheck{$z=0$: $|0-2|=2\ge2$. True. $z=1$: $|1-2|=1\ge2$ is false, correctly
excluded.}

\ans{(a) $y=-9,\;-\tfrac13$ \quad (b) $y=\tfrac72$ \quad
(c) $-4\le z\le 4$ \quad (d) $z\le-2$ or $z=0$ or $z\ge 2$}
\end{solution}

\begin{pitfall}
$|u|\le a$ closes \emph{inwards} to a single interval; $|u|\ge a$ opens
\emph{outwards} to two. Getting these the wrong way round in (c) and (d)
inverts the answer completely.
\end{pitfall}

%---------------------------------------------------------------------
\section{Midpoints, distances and a segment length}

\begin{question}{Q.8, Q.9, Q.10}
\textbf{Q.8} Find the midpoint of the segment joining
(a) $(3,8)$ and $(5,5)$; (b) $(-2,-4)$ and $(2,4)$;
(c) $(6,6)$ and $(-3,-3)$; (d) $(0,9)$ and $(9,0)$.

\smallskip
\textbf{Q.9} Find the distance between
(a) $(12,0)$ and $(0,-8)$; (b) $(4,4)$ and $(-5,-8)$;
(c) $(-2,4)$ and $(1,0)$; (d) $(-7,-2)$ and $(1,-4)$.

\smallskip
\textbf{Q.10} Find the length of the segment joining $C(2,4)$ and $D(4,8)$.
\end{question}

\begin{solution}
\textbf{Q.8 --- midpoints,} $M=\left(\dfrac{x_1+x_2}{2},\dfrac{y_1+y_2}{2}\right)$:

\begin{center}
\small\sffamily
\begin{tabular}{@{}c L{6.4cm} C{3.4cm}@{}}
\toprule
 & Working & Midpoint \\
\midrule
(a) & $\left(\tfrac{3+5}{2},\tfrac{8+5}{2}\right)$      & $(4,\;6.5)$ \\
(b) & $\left(\tfrac{-2+2}{2},\tfrac{-4+4}{2}\right)$    & $(0,\;0)$ \\
(c) & $\left(\tfrac{6-3}{2},\tfrac{6-3}{2}\right)$      & $(1.5,\;1.5)$ \\
(d) & $\left(\tfrac{0+9}{2},\tfrac{9+0}{2}\right)$      & $(4.5,\;4.5)$ \\
\bottomrule
\end{tabular}
\end{center}

\medskip
\textbf{Q.9 --- distances,}
$d=\sqrt{(x_2-x_1)^{2}+(y_2-y_1)^{2}}$:

\begin{center}
\small\sffamily
\begin{tabular}{@{}c L{7.0cm} C{3.6cm}@{}}
\toprule
 & Working & Distance \\
\midrule
(a) & $\sqrt{(0-12)^{2}+(-8-0)^{2}}=\sqrt{144+64}=\sqrt{208}$ & $4\sqrt{13}\approx 14.42$ \\
(b) & $\sqrt{(-5-4)^{2}+(-8-4)^{2}}=\sqrt{81+144}=\sqrt{225}$ & $15$ \\
(c) & $\sqrt{(1+2)^{2}+(0-4)^{2}}=\sqrt{9+16}=\sqrt{25}$      & $5$ \\
(d) & $\sqrt{(1+7)^{2}+(-4+2)^{2}}=\sqrt{64+4}=\sqrt{68}$     & $2\sqrt{17}\approx 8.25$ \\
\bottomrule
\end{tabular}
\end{center}

Parts (b) and (c) are the $9$--$12$--$15$ and $3$--$4$--$5$ right triangles,
which is why they come out whole.

\medskip
\textbf{Q.10.}
\[
|CD|=\sqrt{(4-2)^{2}+(8-4)^{2}}=\sqrt{4+16}=\sqrt{20}=2\sqrt5\approx 4.47 .
\]

\ans{Q.8: $(4,6.5)$, $(0,0)$, $(1.5,1.5)$, $(4.5,4.5)$. \quad
Q.9: $4\sqrt{13}\approx14.42$, $15$, $5$, $2\sqrt{17}\approx8.25$. \quad
Q.10: $2\sqrt5\approx 4.47$.}
\end{solution}

\begin{pitfall}
The distance formula \emph{subtracts} and the midpoint formula \emph{adds}.
Mixing them --- $\left(\tfrac{x_2-x_1}{2},\dots\right)$ for a midpoint --- is
the standard slip, and it silently gives a plausible-looking wrong point.
\end{pitfall}
